%% ============================================================
\section{Real-Time Collaboration Engine}
\label{sec:realtime}
%% ============================================================

\subsection{Overview}

Real-time collaboration in \echoray{} is mediated by a Socket.io server
running within the same Node.js process as the REST API
(\texttt{server/src/server.ts}).  The collaboration engine handles three
categories of events:

\begin{enumerate}[leftmargin=*, label=(\roman*)]
  \item \textbf{Chat messages} --- textual content broadcast to all project
    members in a Socket.io room.
  \item \textbf{AI invocations} --- chat messages prefixed with \texttt{@ai}
    that trigger asynchronous Gemini queries followed by AI-response
    broadcasts.
  \item \textbf{File-tree updates} --- the client persists the current file
    tree via a REST \texttt{PUT} call; in-progress edits flow through the
    Monaco Editor instance without further real-time diffing.
\end{enumerate}

\subsection{Connection Lifecycle and Room Management}

\Cref{alg:socket_lifecycle} presents the connection lifecycle.

\begin{algorithm}[H]
\caption{Socket.io Connection Lifecycle (\texttt{socket.ts})}
\label{alg:socket_lifecycle}
\begin{algorithmic}[1]
\Require{Incoming socket with \texttt{handshake.auth.token} and
         \texttt{handshake.query.projectId}}
\State $token \gets \texttt{socket.handshake.auth.token}$
\State $projectId \gets \texttt{socket.handshake.query.projectId}$
\State $user \gets \textsc{VerifyJWT}(token)$  \Comment{throws on invalid}
\State $project \gets \textsc{GetProjectById}(projectId)$
\If{$user.id \notin project.users$}
  \State \textsc{Disconnect}(socket, \texttt{"unauthorised"})
\EndIf
\State $socket.user \gets user$
\State $socket.project \gets project$
\State $socket.roomId \gets projectId.\textsc{toString}()$
\State $socket.\textsc{join}(socket.roomId)$  \Comment{join project room}
\State \textbf{on} \texttt{"project-message"} $\Rightarrow$ \Call{HandleMessage}{socket, data}
\State \textbf{on} \texttt{"disconnect"} $\Rightarrow$ log departure
\end{algorithmic}
\end{algorithm}

The room model ensures that Socket.io events are scoped to a single project:
the \texttt{socket.broadcast.to(roomId).emit()} call delivers messages to
every socket in the same room except the sender, providing natural isolation
between concurrent projects.

\subsection{Message Handling Protocol}

The \texttt{project-message} event carries a payload conforming to the
following TypeScript interface (defined in \texttt{utils/type.ts}):

\begin{lstlisting}[caption={Message payload type.}, label=lst:msgtype]
interface MessagePayload {
  message: string;          // raw chat text (may contain @ai prefix)
  sender: {
    id: number;
    email: string;
  };
}
\end{lstlisting}

Upon receipt, the server executes the routing logic shown in
\cref{alg:msg_routing}. Non-AI messages are immediately re-broadcast; AI
messages trigger an asynchronous Gemini query before broadcasting.

\begin{algorithm}[H]
\caption{Message Routing (\texttt{socket.ts::HandleMessage})}
\label{alg:msg_routing}
\begin{algorithmic}[1]
\Require{socket with bound room, incoming \texttt{data} payload}
\State $\texttt{socket.broadcast.to}(roomId).\texttt{emit}(\text{"project-message"}, data)$
\If{$data.message.\textsc{startsWith}(\texttt{"@ai"})}$
  \State $prompt \gets data.message.\textsc{replace}(\texttt{"@ai"}, \texttt{""})$
  \State $result \gets \textsc{AwaitAIService.generateResult}(prompt)$
  \State Build $aiPayload$ with $sender = \{id: 0, email: \texttt{"AI"}\}$
  \State $\texttt{io.to}(roomId).\texttt{emit}(\text{"project-message"}, aiPayload)$
\EndIf
\end{algorithmic}
\end{algorithm}

Note the distinction between \texttt{socket.broadcast} (all sockets in room
\emph{except} the sender) and \texttt{io.to} (all sockets in room
\emph{including} the sender). Regular chat messages are rebroadcast with
\texttt{socket.broadcast} because the sender's own UI has already rendered the
message locally; AI responses are emitted with \texttt{io.to} to guarantee the
invoker also receives the response.

\subsection{File-Tree Synchronisation}

The file tree is the central artefact of every project. The synchronisation
strategy in \echoray{} follows a \emph{last-writer-wins} model:

\begin{enumerate}[leftmargin=*, label=\arabic*.]
  \item User edits a file in the Monaco Editor component
    (\texttt{components/CodeEditor.tsx}).
  \item The \texttt{project/[id]/page.tsx} component maintains the entire
    file tree in React state and propagates changes via the
    \texttt{updateFileTree} callback.
  \item On save (explicit user action), the client issues
    \texttt{PUT /api/v1/project/update-file-tree} with the full serialised
    tree.
  \item The server overwrites the \texttt{Project.fileTree} JSON column.
  \item Other collaborators refresh by re-loading the project or receiving
    a future socket event.
\end{enumerate}

This approach sacrifices concurrent-edit conflict resolution (no OT/CRDT) in
favour of implementation clarity. The trade-off is appropriate for the primary
use case---structured pair-programming and code generation sessions---where a
single user typically drives editing while others observe.

\subsection{WebContainer Integration}

Live code preview is provided by the WebContainer API
(\texttt{config/Webcontainer.ts}), which boots a full Node.js~18 runtime
inside the browser's Service Worker context using WebAssembly. The
initialisation is guarded by a singleton pattern to prevent multiple
concurrent boot calls:

\begin{lstlisting}[caption={WebContainer singleton initialisation.},
  label=lst:wc]
let webcontainerInstance: WebContainer | null = null;

export async function getWebContainerInstance(): Promise<WebContainer> {
  if (!webcontainerInstance) {
    webcontainerInstance = await WebContainer.boot();
  }
  return webcontainerInstance;
}
\end{lstlisting}

The file tree—stored in PostgreSQL as the \texttt{Project.fileTree} JSON
object—is directly mountable by calling \texttt{webcontainer.mount(fileTree)},
because the storage format (Equation~\ref{eq:filetree}) matches the
WebContainer \texttt{FileSystemTree} specification verbatim.

\subsection{Socket.io Client Configuration}

On the client side, the Socket.io connection is lazily initialised in
\texttt{config/socketIo.ts}:

\begin{lstlisting}[caption={Socket.io client singleton
  (\texttt{config/socketIo.ts}).}, label=lst:socket_client]
import { io, Socket } from "socket.io-client";

let socket: Socket;

export const initSocket = (
  token: string,
  projectId: string | number
): Socket => {
  socket = io(process.env.NEXT_PUBLIC_API_URL!, {
    auth: { token },
    query: { projectId },
    transports: ["websocket"],
  });
  return socket;
};

export const getSocket = (): Socket => socket;
\end{lstlisting}

Forcing \texttt{transports: ["websocket"]} eliminates the HTTP long-poll
upgrade handshake, reducing connection establishment latency by approximately
one round-trip.

\subsection{Redux State Integration}

The front-end state layer (\texttt{store/}) maintains a minimal Redux slice
(\texttt{userSlice.ts}) tracking the authenticated user's identity. Socket
events update React component state directly via hooks, keeping real-time data
out of Redux to avoid unnecessary re-renders from high-frequency events.
